% --- Frontpage ----------------------------------------------------------------
\maketitle

% --- Prerequisites ------------------------------------------------------------
\chapter*{Prerequisites}
\addcontentsline{toc}{chapter}{Prerequisites}

% --- SECTION ------------------------------------------------------------------
% --- FROM COMPLEX ANALYSIS ----------------------------------------------------
\section*{From Complex Analysis}

\begin{definition*}[$d$-fold map]
    We say a map $f: V \to W$ is a \emph{$d$-fold map} if, for every $w$ in $W$,
    the equation $f(z) = w$ has \emph{exactly} $d$ solutions in $V$ (counting
    multiple solutions by their multiplicity).
\end{definition*}

\begin{theorem*}[Inverse Function Theorem in $\C$]
    \label{thm:inverse_function_theorem_in_C}
    Let $f: U \to \C$ be a holomorphic function on an open set $U \subseteq \C$,
    and let $z_0 \in U$. If $f'(z_0) \neq 0$, then there exist open
    neighborhoods $V \subseteq U$ of $z_0$ and $W$ of $f(z_0)$ such that the
    restriction $f: V \to W$ is a bijection.
    
    Furthermore, the local inverse function $f^{-1}: W \to V$ is holomorphic,
    and its derivative satisfies:
    \begin{equation*}
        (f^{-1})'(f(z_0)) = \frac{1}{f'(z_0)}.
    \end{equation*}
\end{theorem*}

\begin{theorem*}[Maximum Modulus Theorem]
    Let $D \subseteq \C$ be a domain (an open, connected set) and let $f: D \to \C$
    be a holomorphic function. If there exists a point $z_0 \in D$ such that
    $|f(z_0)| \geq |f(z)|$ for all $z \in D$, then $f$ must be a constant
    function on $D$.
\end{theorem*}

\begin{lemma*}[Schwarz's Lemma]
    Let $C = \{ z \in \C \mid |z| < 1 \}$ be the open unit disk, and let
    $f: C \to C$ be a holomorphic function such that $f(0) = 0$. Then, for all
    $z \in C$, we have $|f(z)| \leq |z|$ and $|f'(0)| \leq 1$.

    Furthermore, if $|f(z)| = |z|$ for some non-zero $z \in C$, or if
    $|f'(0)| = 1$, then $f$ is a rotation, meaning $f(z) = \lambda z$ for some
    constant $\lambda \in \C$ with $|\lambda| = 1$.
\end{lemma*}

\begin{theorem*}[Rouché's Theorem]
    Let $D \subseteq \C$ be a domain, and let $\gamma$ be a simple closed contour
    in $D$ whose interior is also entirely contained in $D$. If $f$ and $g$ are
    holomorphic functions on $D$ satisfying
    \[ |f(z) - g(z)| < |f(z)| \]
    for all points $z$ on $\gamma$, then $f$ and $g$ have the exact same number
    of zeros strictly inside $\gamma$, counting multiplicities.
\end{theorem*}

\begin{theorem*}[Hurwitz's Theorem]
    If a sequence of univalent analytic maps $f_n$ on a domain $D$ converges
    uniformly to $f$, then $f$ is either constant or univalent in $D$.
\end{theorem*}

% --- SECTION ------------------------------------------------------------------
% --- FROM TOPOLOGY ------------------------------------------------------------
\section*{From Topology}

To support subsequent global dynamic formulations, we briefly review standard
topological notions adapted to our setting.

% --- SUBSECTION - BASIC CONCEPTS ----------------------------------------------
\subsection*{Basic concepts}

\begin{definition*}[Topology]
    Let $X$ be a set. A \emph{topology} on $X$ is a collection $\tau$ of subsets
    of $X$ that satisfies the following three properties:
    \begin{itemize}
        \item The empty set $\emptyset$ and the set $X$ itself belong to $\tau$.
        \item The union of any arbitrary collection of sets in $\tau$ also
        belongs to $\tau$.
        \item The intersection of any finite collection of sets in $\tau$ also
        belongs to $\tau$.
    \end{itemize}
\end{definition*}

\begin{definition*}[Topological space]
    A \emph{topological space} is a pair $(X, \tau)$ consisting of a set $X$ and
    a topology $\tau$ on $X$. When the topology $\tau$ is clear from the
    context, we often simply refer to the space $X$.
\end{definition*}

\begin{definition*}[Open set]
    Let $(X, \tau)$ be a topological space. A subset $U \subseteq X$ is called
    an \emph{open set} if it belongs to the topology $\tau$, that is, if $U \in
    \tau$.
\end{definition*}

\begin{definition*}[Closed set]
    A subset $K \subseteq X$ of a topological space is called a \emph{closed
    set} if its complement, $X - K$, is an open set.
\end{definition*}

\begin{definition*}[Closure]
    Let $A$ be a subset of a topological space $X$. The \emph{closure} of $A$,
    denoted by $\closure{A}$, is the intersection of all closed sets containing
    $A$. Equivalently, it is the smallest closed set containing $A$.
\end{definition*}

\begin{definition*}[Dense set]
    A subset $A$ of a topological space $X$ is called \emph{dense} in $X$ if its
    closure is the entire space, that is, if $\closure{A} = X$.
\end{definition*}

\begin{definition*}[Interior]
    Let $A$ be a subset of a topological space $X$. The \emph{interior} of $A$,
    denoted by $\interior{A}$, is the union of all open
    sets contained within $A$. Equivalently, it is the largest open set contained
    in $A$.
\end{definition*}

\begin{definition*}[Boundary]
    Let $A$ be a subset of a topological space $X$. The \emph{boundary} (or
    frontier) of $A$, denoted by $\boundary{A}$, is the set of points that are in
    the closure of $A$ but not in the interior of $A$. Thus, $\boundary{A} =
    \closure{A} - \interior{A}$.
\end{definition*}

\begin{definition*}[Accumulation points and derived set]
    Let $A \subseteq X$. A point $x \in X$ is an \emph{accumulation point} (or
    limit point) of $A$ if every neighborhood of $x$ contains at least one point
    of $A$ distinct from $x$. The set of all accumulation points of $A$ is
    called the \emph{derived set} of $A$ and is denoted by $\derived{A}$.
\end{definition*}

\begin{definition*}[Isolated points]
    Let $A \subseteq X$. A point $x \in A$ is an \emph{isolated point} of $A$ if
    it is not an accumulation point of $A$; equivalently, if there exists a
    neighborhood of $x$ that contains no other points of $A$ besides $x$.
\end{definition*}

\begin{definition*}[Neighborhood]
    A subset $N$ of a topological space $X$ is a \emph{neighborhood} of a point
    $x \in X$ if there exists an open set $U$ such that $x \in U \subseteq N$.
    If $N$ itself is an open set, it is frequently referred to as an \emph{open
    neighborhood}.
\end{definition*}

% --- SUBSECTION - METRIC SPACES -----------------------------------------------
\subsection*{Metric spaces}

\begin{definition*}[Metric or distance]
    Let $X$ be a non-empty set. A \emph{metric} (or \emph{distance}) (on $X$) is
    a mapping
    \begin{alignat*}{2}
        \sigma: X \times X &{}\to{}     &&[0, \infty)  \\
                    (x, y) &{}\mapsto{} &&\sigma(x, y)
    \end{alignat*}
    verifying the following for every $x, y, z \in X$:
    \begin{itemize}
        \item $\sigma(x, y) = 0 \iff x = y$;
        \item $\sigma(x, y) = \sigma(y,x)$; and
        \item $\sigma(x, z) \leq \sigma(x, y) + \sigma(y, z)$.
    \end{itemize}
\end{definition*}

\begin{definition*}[Metric space]
    A metric space is a tuple $(X, \sigma)$ where $X$ is a non-empty set and
    $\sigma$ is a metric on $X$.
\end{definition*}

\begin{definition*}[Distance from a point to a set]
    Let $(X, \sigma)$ be a metric space, $A \subseteq X$ a non-empty set, and $x
    \in X$. The \emph{distance} between the point $x$ and the set $A$ is defined
    as
    \begin{equation*}
        \sigma(x, A) = \inf_{y \in A} \sigma(x, y).
    \end{equation*}
\end{definition*}

\begin{definition*}[Distance between two sets]
    Let $A$ and $B$ be non-empty subsets of a metric space $(X, \sigma)$. The
    \emph{distance} between $A$ and $B$ is defined as
    \begin{equation*}
        \sigma(A, B) = \inf_{\substack{x \in A \\ y \in B}} \sigma(x, y).
    \end{equation*}
\end{definition*}

\begin{definition*}[Diameter of a set]
    Let $A$ be a non-empty subset of a metric space $(X, \sigma)$. The
    \emph{diameter} of $A$ is defined as
    \begin{equation*}
        \diameter(A) = \sup_{x, y \in A} \sigma(x, y).
    \end{equation*}
\end{definition*}

\begin{definition*}[Ball]
    Let $(X, \sigma)$ be a metric space, $x \in X$, and $r > 0$. The \emph{open
    ball} (or simply the \emph{ball}) of radius $r$ centered at $x$ is the set
    \begin{equation*}
        B(x, r) = \{ y \in X \mid \sigma(x, y) < r \}.
    \end{equation*}
    Similarly, the \emph{closed ball} of radius $r$ centered at $x$ is the set
    \begin{equation*}
        \closure{B}(x, r) = \{ y \in X \mid \sigma(x, y) \leq r \}.
    \end{equation*}
\end{definition*}

\begin{definition*}[Bounded set]
    A subset $A$ of a metric space $(X, \sigma)$ is \emph{bounded} if its
    diameter is finite, that is, if $\diameter(A) < \infty$. Equivalently, $A$
    is bounded if it is contained in a ball (open or closed).
\end{definition*}

\begin{definition*}[Open set]
    A subset $U$ of a metric space $(X, \sigma)$ is \emph{open} if, for every
    point $x \in U$, there exists a radius $r > 0$ such that the ball $B(x, r)$
    is entirely contained in $U$.
\end{definition*}

% --- SUBSECTION - COMPACTNESS -------------------------------------------------
\subsection*{Compactness}

\begin{definition*}[Compact Set]
    A set $K$ is \emph{compact} if every open cover of $K$ contains a finite
    subcover. In a metric space, this is equivalent to \emph{sequential
    compactness}, where every sequence in $K$ has a convergent subsequence whose
    limit is in $K$.
\end{definition*}

By the Heine-Borel theorem, compactness for a subset $K \subseteq \C$ is
equivalent to $K$ being closed and bounded. However, because the extended plane
$\ExtC$ is itself already compact (being homeomorphic to the sphere $S^2$),
bounds are redundant. The transformation from $\C$ to $\ExtC$ represents a
standard one-point or Alexandroff compactification.

\begin{definition*}[Compact set in $\ExtC$]
    A subset $K \subseteq \ExtC$ is compact if and only if it is closed in the
    chordal topology.
\end{definition*}

% --- SUBSECTION - CONNECTEDNESS -----------------------------------------------
\subsection*{Connectedness}

\begin{definition*}[Connected Set]
    A set $A$ is \emph{connected} if it cannot be partitioned by two disjoint,
    non-empty open sets $V$ and $W$ such that $A \subseteq V \cup W$.
\end{definition*}

\begin{definition*}[Component of a set]
    Given a set $B$, we say a subset $A$ of $B$ is a component of $B$ if it is a
    maximal connected subset.
\end{definition*}

\begin{definition*}[Path-Connected Set]
    A set $A \subseteq \ExtC$ is \emph{path-connected} if for every pair of
    points $z_0, z_1 \in A$, there exists a continuous map (a path) $\gamma: [0,
    1] \to A$ such that $\gamma(0) = z_0$ and $\gamma(1) = z_1$.
\end{definition*}

In general topology, path-connectedness is a strictly stronger condition than
connectedness; every path-connected set is connected, but the converse is false
(as exemplified by the Topologist's Sine Curve).

\begin{definition*}[Domain]
    An open, connected subset $D \subseteq \ExtC$ is called a \emph{domain}.
\end{definition*}

The regular geometry of open sets yields an equivalence vital to complex
dynamics.

\begin{theorem*}[An open set is connected $\iff$ is path-connected]
    Let $U \subseteq \ExtC$ be an open set. Then $U$ is connected if and only if
    it is path-connected.
\end{theorem*}

Consequently, any domain $D$ is automatically path-connected. Thus, domains
could also be defined as open, path-connected subsets of $\ExtC$.

% --- SUBSECTION - JORDAN CURVES -----------------------------------------------
\subsection*{Jordan Curves}

To properly interpret the connectivity of domains and properties of the Fatou
and Julia sets, we recall the basic topology of curves on the Riemann sphere
$\ExtC$.

\begin{definition*}[Jordan Arc and Jordan Curve]
    A \emph{Jordan arc} is the image of a continuous, injective map $\gamma: [0,
    1] \to \ExtC$. A \emph{Jordan curve} (or a simple closed curve) is the image
    of a continuous map $\gamma: [0, 1] \to \ExtC$ that is injective on $[0, 1)$
    and satisfies $\gamma(0) = \gamma(1)$.
\end{definition*}

The fundamental global property of Jordan curves is their ability to divide the
sphere into distinct regions, a result which is heavily utilized in complex
dynamics when defining domain connectivity.

\begin{theorem*}[Jordan Curve Theorem]
    \label{thm:jordan_curve_theorem}
    Let $\gamma$ be a Jordan curve in $\ExtC$. Then its complement $\ExtC -
    \gamma$ consists of exactly two disjoint connected components, each of which
    is an open domain having $\gamma$ as its common boundary.
\end{theorem*}

While the Jordan Curve Theorem guarantees exactly two components, complex
dynamics frequently requires these components to be topologically well-behaved.
This regular geometry is guaranteed by a stronger two-dimensional extension.

\begin{theorem*}[Schoenflies Theorem for the Sphere]
    \label{thm:schoenflies_theorem}
    Let $\gamma$ be a Jordan curve in $\ExtC$. Then there exists a homeomorphism
    $\Phi: \ExtC \to \ExtC$ that maps $\gamma$ directly onto the standard
    equatorial circle. Consequently, the two components of $\ExtC - \gamma$
    are homeomorphic to open discs, and their closures are homeomorphic to closed
    discs (topological discs).
\end{theorem*}

% --- Table of Contents --------------------------------------------------------
\tableofcontents